\section{Discussion}\label{sec:discussion}

Sections~\ref{sec:abm} and~\ref{sec:hazard} presented the paper's central contrast. Household simulation, rational or behaviorally extended, does not close the gap to the empirical benchmark. A structurally distinct hazard framework---no household utility, a substantially larger non-synthetic sample---recovers 91.3\% of it in microsimulation form on the benchmark-consistent accounting basis at the headline off-window floor, and 97.9\% at the in-sample calibration point. Across that floor's range the recovery runs 88.2--93.7\%, a miss of 6.3 to 11.8 points. But no estimator's monthly path co-movement, Path B's included, survives removing the window's shared trend (Section~\ref{sec:pathb}). The contrast is easy to over-read. The $\beta_1 = 0$ null recovers 85.7\% on the same basis at that headline floor, 88.7\% in-sample. Components common to any prepayment model---scheduled amortization, the seasoning baseline, the turnover floor---therefore deliver the large majority of the level, and the caps sat far above what any plausible prepayment environment would have produced. Named loan-level mechanisms are measured small wherever I can ablate them. Switching burnout off moves Path B by \$6.58 billion ($0.86$ points) at the in-sample floor and \$3.48 billion ($0.45$ points) at the headline off-window floor. Zeroing the FICO/LTV covariates moves it $1.3$ points. The delinquency transition matrix is bounded at \$0.39 billion (Appendix~\ref{app:params}), and no servicer channel beyond that pipeline is modeled at all. Three mechanisms therefore remain \emph{candidates} for the residual and for the timing failure every estimator shares, not demonstrated drivers of the level. That residual is modest at the in-sample calibration (2.1 points) but 8.7 points at the headline off-window calibration, roughly 1.6$\times$ the headline marginal. Servicer behavior beyond the delinquency pipeline, updated-equity dynamics, and home-equity extraction carry no bound of any kind here, each plausibly rate-cycle-correlated. First, then, loan-level heterogeneity in credit score, geography, and life circumstance makes mobility no smooth function of coupon and market rate alone; the hazard framework's 295 stratum fixed effects are one mechanism letting Path A track pool-level heterogeneity the ABM's coupon-cohort buckets cannot represent. Second, servicer behavior---loss mitigation, forbearance, and modification practices that vary by servicer and were unusually active in the higher-rate, higher-delinquency-risk QT window---reaches prepayment and default outcomes outside any borrower-side decision model, which Path B's competing-risks structure and Markov delinquency pipeline represent directly if partially. Third, prepayment burnout is representable at the cohort level, as in both hazard paths, but not inside a closed, per-agent household population, as Section~\ref{sec:abm-results}'s falsified survivor-selection experiment showed.

\subsection{The Dual Systemic Fallout}\label{sec:discussion-fallout}

The macroeconomic literature often focuses on the housing mobility implications of the lock-in effect. But its transmission into the fixed-income sector creates a two-pronged systemic risk.

The first channel runs through the central bank itself. Extending MBS durations deprives the SOMA of its scheduled cash flows, and it breaks the predictability effective Quantitative Tightening requires; Appendix~\ref{sec:robustness-wal} states the extension in approximate weighted-average-life units. The \$764.7 billion shortfall-versus-cap is, per the decomposition above, mostly rate-inelastic---a gap to a cap schedule no plausible prepayment environment would have met. The rate-driven object is the duration extension itself: a 9.4-year portfolio WAL under the empirical path against 3.4 years at 2021 realized speeds (Appendix~\ref{sec:robustness-wal}). That 3.4-year comparator is a refinancing-boom baseline. The extension, not the cap arithmetic, feeds the fragility channels below.

The second channel runs through private banks, which, unlike the Fed, must mark to market: the structural duration extension of these low-coupon MBS is one source of the unrecognized mark-to-market losses that accumulated across the banking sector. \citet{jiang2023} show duration extension alone was not sufficient to precipitate failure---it was mark-to-market losses interacting with a high share of uninsured deposits that made specific regional institutions vulnerable to self-fulfilling runs, uninsured depositors withdrawing once reported asset values no longer covered liabilities. The duration extension that is a necessary precondition for this fragility channel comes from the below-market coupon stock meeting rate normalization under any payoff rule: Section~\ref{sec:abm-danish}'s rule-only counterfactual shortens the portfolio by only 0.6 of the roughly six extension-years (Appendix~\ref{sec:robustness-wal}), and those six years are measured against 2021's refinancing-boom speeds, not normal turnover. Read instead against the baseline turnover floor, the empirical path's approximate weighted-average life is 9.4 years against 9.5 at the window open and 8.5 against 8.6 at its close, so the extension measures refinancing-versus-turnover, not lock-in-versus-normalcy. The elasticity channel I identify is therefore a second-order contributor to that precondition, and neither it nor the extension is, on its own, a sufficient condition for bank failure. On the institutional side I price cash-flow timing alone (principal arriving later than the cap schedule assumed), not a market-value loss. Appendix~\ref{sec:robustness-wal} declines option-adjusted duration and convexity as out of scope, so no mark-to-market or DV01 image is put on that 0.6 of an extension-year. The bank-side losses \citet{jiang2023} price have no counterpart here. The fragility channel is invoked as context, not as something this design measures.

Beyond fixed-income transmission, the lock-in effect imposes a real economic cost by keeping households from relocating toward higher-productivity labor markets. \citet{hsieh2019} estimate substantial aggregate output losses from such spatial misallocation of labor under housing supply constraints in high-productivity U.S. cities. Their headline magnitudes are cumulative 1964--2009 level effects revised in a published correction, so I cite the mechanism, not a point estimate. Theirs is a supply-side housing-availability constraint; the demand-side mobility constraint I quantify runs through the same channel, since households who cannot economically justify relocating cannot respond to geographic labor-demand differences, whatever the housing supply elsewhere. No per-household or per-move welfare estimate enters the paper, so the labor-reallocation leg of the mobility argument is directional only. The per-move surplus that would turn this leg into a magnitude is nowhere quantified. Nor do I translate the \$764.7 billion shortfall into welfare terms: it is an accounting gap against a cap schedule, mostly rate-inelastic by the decomposition above, and neither a lower nor an upper bound on the par-payoff rule's welfare cost. The welfare-relevant objects are the constrained moves, whose fiscal image is the lock-in marginal ($+\$42.6$ billion of trapped roll-off at the headline off-window floor, $+\$70.3$ billion at the in-sample calibration point), and the forgone labor-reallocation gains this literature identifies, which my accounting framework does not measure.

A back-of-envelope prices the extension in fiscal units. At the window's policy-rate levels the funding spread over the book's 2.49\% weighted-average coupon ran roughly 250--300 basis points. The \$42.6 billion of trapped roll-off therefore carries on the order of \$1 billion per year of negative carry on the consolidated public balance sheet, for as long as the gap persists. That figure restates in remittances terms the maturity extension Table~\ref{tab:wal} prices in years. It derives from committed quantities---the book coupon, the marginal---and the public policy-rate path. It is subject to the same face-versus-cash incidence reading as the buyback bracket (Appendix~\ref{app:params}).
